%%% FIELD thm-chen-proposed-statement
The full heterogeneous finite-observed certified class \(\mathcal M_{\mathrm{cert}}^{(q)}(K,J)\) and \(\mathcal M_{\mathrm{Chen}}\) are incomparable in the ambient mixed-type latent-structure universe.  The noncontainment \(\mathcal M_{\mathrm{cert}}^{(q)}(K,J)\not\subseteq\mathcal M_{\mathrm{Chen}}\) is witnessed by the equal-alphabet dense ternary member: in the equal-alphabet subfamily every selected observed variable and every latent coordinate have support cardinality \(q\), so Chen et al.'s discrete sufficient-measurement requirement \(\lvert\operatorname{supp}(X_j)\rvert>\lvert\operatorname{supp}(H_k)\rvert\) fails, and the ternary witness is outside their two-sufficient-measurement class.  Conversely, \(\mathcal M_{\mathrm{Chen}}\) contains members with continuous measurements and permits graphs without two disjoint triangular views, whereas \(\mathcal M_{\mathrm{cert}}^{(q)}(K,J)\) requires finite observed alphabets.  Completeness in Chen et al. supplies a discretization with full-rank tables, whereas this project derives full rank from the finite local certificate \(\Delta_{\mathrm{DT}}\) without assuming completeness or larger-support measurements.
%%% FIELD thm-chen-statement-replace
The full heterogeneous finite-observed certified class \(\mathcal M_{\mathrm{cert}}^{(q)}(K,J)\) and \(\mathcal M_{\mathrm{Chen}}\) are incomparable in the ambient mixed-type latent-structure universe.  The noncontainment \(\mathcal M_{\mathrm{cert}}^{(q)}(K,J)\not\subseteq\mathcal M_{\mathrm{Chen}}\) is witnessed by the equal-alphabet dense ternary member: in the equal-alphabet subfamily every selected observed variable and every latent coordinate have support cardinality \(q\), so Chen et al.'s discrete sufficient-measurement requirement \(\lvert\operatorname{supp}(X_j)\rvert>\lvert\operatorname{supp}(H_k)\rvert\) fails, and the ternary witness is outside their two-sufficient-measurement class.  Conversely, \(\mathcal M_{\mathrm{Chen}}\) contains members with continuous measurements and permits graphs without two disjoint triangular views, whereas \(\mathcal M_{\mathrm{cert}}^{(q)}(K,J)\) requires finite observed alphabets.  Completeness in Chen et al. supplies a discretization with full-rank tables, whereas this project derives full rank from the finite local certificate \(\Delta_{\mathrm{DT}}\) without assuming completeness or larger-support measurements.
%%% FIELD proof-thm-chen-class-relation
Write
\[
 \mathcal C_{\mathrm{eq}}^{(q)}(K,J)
 =\bigl\{(G,P)\in\mathcal M_{\mathrm{cert}}^{(q)}(K,J):m_j=q
   \text{ for every }j\bigr\}.
\]
This is a subfamily of the full heterogeneous class, not a replacement for it.

We first prove \(\mathcal M_{\mathrm{cert}}^{(q)}(K,J)\not\subseteq\mathcal M_{\mathrm{Chen}}\).  Proposition~\textnormal{prop:ternary-witness} gives a rational faithful law \(\widetilde P\), on the fixed graph \(\Gamma^{\star}\), in
\(\mathcal C_{\mathrm{eq}}^{(3)}(3,8)\).  In particular,
\[
 (\Gamma^{\star},\widetilde P)
 \in\mathcal C_{\mathrm{eq}}^{(3)}(3,8)
 \subseteq\mathcal M_{\mathrm{cert}}^{(3)}(3,8).
\]
Strictly positive latent probabilities give
\(\lvert\operatorname{supp}_{\widetilde P}(H_k)\rvert=3\) for every \(k\).  Every observed alphabet in this witness is \(\{0,1,2\}\), and the chosen neighborhood retains strictly positive measurement cells, so
\(\lvert\operatorname{supp}_{\widetilde P}(X_j)\rvert=3\) for every \(j\).  By Lemma~\textnormal{lem:chen-mixed-sufficient-measurement}, a discrete sufficient measurement for a latent parent must have strictly larger support.  Here that inequality would be \(3>3\), which is false.  Thus none of the witness measurements is sufficient in Chen et al.'s discrete sense, the two-sufficient-measurement requirement fails, and
\[
 (\Gamma^{\star},\widetilde P)\notin\mathcal M_{\mathrm{Chen}}.
\]
Since the witness belongs to the full certified class, this single equal-alphabet member proves noncontainment of the full class.  Notice that no support-cardinality claim has been made about certified models with \(m_j>q\).

For the reverse noncontainment, it is useful to exhibit a member of the continuous branch recorded in Lemma~\textnormal{lem:chen-mixed-sufficient-measurement}.  Let \(H\) be uniform on \(\mathcal Q\), choose distinct real numbers \(\mu_0,\ldots,\mu_{q-1}\), and, conditionally on \(H=h\), let three pure children be independent with
\[
 X_i=\mu_h+\epsilon_i,\qquad
 \epsilon_1,\epsilon_2,\epsilon_3\stackrel{\mathrm{iid}}{\sim}N(0,1).
\]
Relabel the states so that \(\mu_0<\cdots<\mu_{q-1}\).  The conditional densities \(f_h(x)=\phi(x-\mu_h)\) are linearly independent.  Indeed, if
\(\sum_h a_h f_h(x)=0\) almost everywhere, continuity makes the equality hold for every \(x\).  Multiplication by \(\sqrt{2\pi}\exp(x^2/2)\) gives
\[
 \sum_{h=0}^{q-1}a_h\exp(-\mu_h^2/2)\exp(\mu_h x)=0.
\]
Divide by \(\exp(\mu_{q-1}x)\) and let \(x\to\infty\); every term but the last tends to zero, so \(a_{q-1}=0\).  Repetition proves \(a_h=0\) for every \(h\).  Consequently, if \(g:\mathcal Q\to\mathbb R\) satisfies \(E\{g(H)\mid X_i\}=0\), multiplication by the positive marginal density of \(X_i\) yields
\[
 \sum_{h=0}^{q-1}q^{-1}g(h)f_h(x)=0
 \quad\text{for almost every }x,
\]
and linear independence gives \(g=0\).  Thus each continuous child satisfies the required completeness condition and has distinct conditional laws across latent states; in particular, any two of the three children are sufficient measurements in the continuous branch.

This star law is also faithful.  Given any subset \(C\) of the other children and any value \(x_C\), the posterior weights are
\[
 w_h(x_C)
 =\frac{q^{-1}\prod_{i\in C}\phi(x_i-\mu_h)}
        {\sum_{r=0}^{q-1}q^{-1}\prod_{i\in C}\phi(x_i-\mu_r)}>0.
\]
Because the \(\mu_h\) are distinct, \(\operatorname{Var}_{w(x_C)}(\mu_H)>0\).  For two children not in \(C\), conditional independence given \(H\) and the zero-mean independent noises imply
\[
 \operatorname{Cov}(X_i,X_j\mid X_C=x_C)
 =\operatorname{Var}_{w(x_C)}(\mu_H)>0.
\]
Also, the conditional law \(N(\mu_h,1)\) of any remaining child depends on \(h\), while every posterior weight is positive, so \(H\) and that child are dependent given \(X_C=x_C\).  These two calculations rule out every conditional independence not implied by d-separation in a one-parent star; conditioning on \(H\) gives precisely the child independences implied by the Markov factorization.  The model is therefore faithful.  It has a one-factor latent set and three pure children, so it meets the support-identification branch as well as the completeness, nondegeneracy, and two-sufficient-measurement conditions defining \(\mathcal M_{\mathrm{Chen}}\).  Hence it is a member of \(\mathcal M_{\mathrm{Chen}}\).

Each \(X_i\) in this member has support \(\mathbb R\).  By Definition~\textnormal{def:certified-class} and the underlying finite-observed model class, every observed coordinate of \(\mathcal M_{\mathrm{cert}}^{(q)}(K,J)\) instead has a finite state space \(\mathcal X_j=\{0,\ldots,m_j-1\}\).  The displayed continuous member therefore belongs to no such certified class, and
\[
 \mathcal M_{\mathrm{Chen}}\not\subseteq
 \mathcal M_{\mathrm{cert}}^{(q)}(K,J).
\]
The cited Chen lemma also records that their class does not impose two disjoint triangular views, providing a separate structural distinction.

Finally, Lemma~\textnormal{lem:chen-mixed-sufficient-measurement} attributes Chen et al.'s full-rank discretization to completeness.  Here Theorem~\textnormal{thm:local-factorization} derives rectangular view rank from the local Gram minors, and Theorem~\textnormal{thm:certificate-frontier} packages their product as \(\Delta_{\mathrm{DT}}\), with necessity only in the equal-alphabet case.  Thus the two rank mechanisms are different, and both asserted noncontainments follow.
%%% FIELD proof-prop-ternary-witness
The rows \(1,2,3\) of \(\Gamma^{\star}\) have diagonal parents \(1,2,3\), respectively, in the order \(\rho_1=(1,2,3)\), and none has a parent later than its diagonal parent.  The rows \(4,5,6\) have diagonal parents \(2,1,3\), respectively, in the order \(\rho_2=(2,1,3)\), again with no later parent.  Thus \(S_1,S_2\) satisfy double triangularity.  Rows \(7,8\) have parent sets \(\{1,3\}\) and \(\{2,3\}\), so \(S_3\) covers every latent coordinate.

Directly from \(\Gamma^{\star}\), the observed-child sets of \(H_1,H_2,H_3\) are
\[
 C_1=\{1,2,3,5,6,7\},\qquad
 C_2=\{2,3,4,5,6,8\},\qquad
 C_3=\{3,6,7,8\}.
\]
For \(C_1,C_2\), the elements \(1\) and \(4\) witness the two failed containments; for \(C_1,C_3\), the elements \(1\) and \(8\) do so; and for \(C_2,C_3\), the elements \(4\) and \(7\) do so.  Hence the subset condition holds.  Every child of \(H_3\), namely \(3,6,7,8\), has at least one additional latent parent, so \(H_3\) has no pure child.

Fix an observed row \(j\).  For every parent \(\ell\),
\[
 \sum_{x=0}^2\left(\mathbf 1\{x=h_\ell\}-\frac13\right)=0,
\]
so the displayed formula for \(P^{\star}(X_j=x\mid H=h)\) sums to one.  Put
\(S_j=\sum_{\ell:\Gamma^{\star}_{j\ell}=1}c^{\star}_{j\ell}\).  The geometric definition of the coefficients gives
\[
 S_j\le 2^{-(4+j)}\sum_{\ell=1}^3 2^{-3\ell}<\frac1{224}.
\]
Since \(-1/3\le \mathbf 1\{x=h_\ell\}-1/3\le2/3\), every cell obeys
\[
 \frac13-\frac{S_j}{3}
 \le P^{\star}(X_j=x\mid H=h)
 \le\frac13+\frac{2S_j}{3}.
\]
The lower bound is positive and the upper bound is less than one.  Together with \(\pi_h^{\star}=3^{-3}>0\), the empty latent DAG, and conditional independence of the observed coordinates given \(H\), this proves that \(P^{\star}\) is a strictly positive law with the required Markov factorization.

For measurement nondegeneracy, take two different configurations of the parents of row \(j\), and let \(\ell_0\) be their smallest differing parent index.  Evaluate the difference of their conditional mass functions at the state equal to the first configuration's value of \(h_{\ell_0}\).  All contributions with \(\ell<\ell_0\) cancel, while the \(\ell_0\) contribution has absolute value \(c^{\star}_{j\ell_0}\).  The absolute value of the sum of all later contributions is at most
\[
 \sum_{\substack{\ell>\ell_0\\\Gamma^{\star}_{j\ell}=1}}c^{\star}_{j\ell}
 \le c^{\star}_{j\ell_0}\sum_{r=1}^{\infty}2^{-3r}
 =\frac{c^{\star}_{j\ell_0}}7
 <c^{\star}_{j\ell_0}.
\]
The two conditional mass functions therefore differ.

For a selected row, double triangularity says that all its parents other than its diagonal parent occur in the predecessor context.  Separating the diagonal term in the formula for \(P^{\star}\) consequently gives, for every displayed context,
\[
 B_{a,k,u}
 =b^{\star}_{s(a,k),u}\mathbf 1^{\mathsf T}
  +c^{\star}_{s(a,k),\rho_a(k)}Q_3.
\]
For any probability vector \(b\) and scalar \(t\),
\[
 b\mathbf 1^{\mathsf T}+tQ_3
 =tI_3+\left(b-\frac t3\mathbf 1\right)\mathbf 1^{\mathsf T}.
\]
When \(t>0\), multilinearity of the determinant in its columns shows directly that
\(\det(tI_3+u\mathbf 1^{\mathsf T})=t^3+t^2\mathbf 1^{\mathsf T}u\): the term using no column of the rank-one matrix is \(t^3\), the three terms using one such column sum to \(t^2\mathbf 1^{\mathsf T}u\), and every term using two is zero because those columns are proportional.  Taking \(u=b-t\mathbf 1/3\) and using \(\mathbf 1^{\mathsf T}b=1\) yields
\[
 \det\!\left(b\mathbf 1^{\mathsf T}+tQ_3\right)
 =t^3\left(1+t^{-1}\mathbf 1^{\mathsf T}
       \left(b-\frac t3\mathbf 1\right)\right)
 =t^2.
\]
Therefore
\[
 \det(B_{a,k,u})
 =\bigl(c^{\star}_{s(a,k),\rho_a(k)}\bigr)^2>0.
\]
Definition~\textnormal{def:local-minors} gives \(c_{a,k,u}=\det(B_{a,k,u})^2>0\), and Definition~\textnormal{def:certificate} then gives \(\Delta_{\mathrm{DT}}>0\).

Work now in the finite-dimensional product of conditional-table simplices for the fixed graph.  Strict cell positivity, inequality of each of the finitely many pairs of measurement columns, and nonvanishing of the finitely many displayed determinants are relatively open conditions.  The graph conditions are fixed.  Hence all the support, graph, nondegeneracy, and certificate properties just verified hold throughout some relative neighborhood \(U\) of the table parameter of \(P^{\star}\).

By Lemma~\textnormal{lem:positive-faithful-parameters-dense}, faithful strictly positive parameters are relatively open and dense, so \(U\) contains a nonempty relatively open set \(U_{\mathrm f}\) of faithful parameters.  The affine simplex coordinates and all defining equations have rational coefficients; rational points are dense in their relative interiors.  Thus \(U_{\mathrm f}\) contains a rational table parameter.  Its induced joint probabilities are finite sums and products of rational numbers and are therefore rational.  Every such chosen law is faithful and retains every property defining \(U\); by Definition~\textnormal{def:certified-class} it belongs to \(\mathcal M_{\mathrm{cert}}^{(3)}(3,8)\).  Theorem~\textnormal{thm:fixed-k-identification} identifies its recovery target.  Since the graph still has no pure child of \(H_3\), this is the required dense ternary no-pure-child witness.  The argument asserts only that every neighborhood contains a rational faithful perturbation, not that the displayed parameter \(P^{\star}\) itself is faithful.
%%% FIELD proof-lem-positive-faithful-parameters-dense
Let \(\Omega\) be the affine product of the open conditional-table simplices for the fixed DAG \(G\).  In each table omit its last cell, so that \(\Omega\) is an open subset of a finite-dimensional real affine space with coordinates \(\theta\).  By the Markov factorization in Assumption~\textnormal{ass:markov}, each joint cell probability is a polynomial in \(\theta\), after the omitted cells are replaced by one minus the sum of the retained cells.

There are only finitely many triples \((V,W,C)\), where \(V,W\) are distinct vertices and \(C\) is a subset of the remaining vertices.  For fixed states \(v,v',w,w',c\), define
\[
 F_{v,v',w,w',c}(\theta)
 =P_\theta(v,w,c)P_\theta(v',w',c)
  -P_\theta(v,w',c)P_\theta(v',w,c),
\]
where variables outside \(\{V,W\}\cup C\) are summed out.  This is a polynomial.  Because all cells are positive on \(\Omega\),
\[
 V\mathrel{\perp}_{P_\theta}W\mid C
 \quad\Longleftrightarrow\quad
 F_{v,v',w,w',c}(\theta)=0
 \text{ for every }v,v',w,w',c.
\]

We next show directly that, whenever \(V\) and \(W\) are d-connected given \(C\), at least one polynomial in this collection is nonzero.  Choose a simple active path and split it at its colliders into collider-free sections.  A collider-free section is a trek: it has a source from which its arrows point along the two arms.  At a boundary point of the table polytope, give that source a fair bit and make every other vertex on the section copy the bit along the relevant arrow.  At each collider, set its bit equal to the exclusive-or of the two adjacent section bits.  If the collider is not itself in \(C\), activity supplies a directed path from it to a descendant in \(C\); transmit the collider bit along such a path.  Where several transmission paths merge, transmit their exclusive-or, processing vertices in topological order.  Assign each collider to one chosen endpoint in \(C\), and set every unused vertex and every unused member of \(C\) to zero.  These assignments are deterministic functions of parents and hence define a boundary Markov table for \(G\).

Conditioning the members of \(C\) to be zero imposes, group by group, that the exclusive-or of the assigned collider bits is zero.  Taking the exclusive-or of all these equations makes every internal section bit occur twice and cancel; the result is equality of the endpoint bits at \(V\) and \(W\).  The conditioning event has positive probability.  Moreover, simultaneously flipping all section-source bits leaves every collider bit and every conditioning equation unchanged, so both endpoint values occur with positive conditional probability.  Hence the conditional \(2\)-by-\(2\) table of \((V,W)\) has positive diagonal entries and zero off-diagonal entries, and one of its cross-product polynomials is nonzero.  For \(q>2\), use states \(0,1\) in this construction and give the remaining states probability zero at the boundary.  Thus the same polynomial, regarded as a polynomial on the \(q\)-state parameter space, is not identically zero.  This establishes the required active-path witness for every d-connected triple.

For each d-connected triple choose one such nonzero polynomial \(F_{V,W,C}\).  A parameter that is Markov but unfaithful must satisfy an additional conditional independence for some d-connected triple and therefore lies in
\[
 \bigcup_{(V,W,C):\,V\not\!\perp_G W\mid C}
 \{\theta\in\Omega:F_{V,W,C}(\theta)=0\}.
\]
This is a finite union.  A nonzero real polynomial on \(\mathbb R^d\) has a zero set of Lebesgue measure zero: for \(d=1\) it has finitely many roots, and induction on \(d\), after expanding in the last coordinate and applying Fubini outside the common zero set of its nonzero coefficient polynomials, proves the claim.  The same induction shows that its zero set has empty interior.  Restricting to the affine coordinates of \(\Omega\) preserves both conclusions.  Therefore the unfaithful parameters have measure zero and empty relative interior, and faithful parameters are dense and have full relative Lebesgue measure.

Finally, for a fixed triple, failure of its conditional independence means that at least one of finitely many continuous polynomials is nonzero, an open condition.  Faithfulness is the intersection of these conditions over finitely many d-connected triples, so the faithful locus is relatively open.  The Markov factorization already supplies every d-separation independence; exclusion of all additional independences therefore proves precisely faithfulness to \(G\).
